\begin{theorem}[The real three-manifolds]\hypertarget{tag-02PG}{}\label{mapsI:src:A.3.3}
\StatementRecord{mapsI:src:A.3.3}{proved}{02PG}{\SeriesRef{mapsI:src:A.3.2}}

Every nonempty fixed set in Theorem~\ref{mapsI:src:A.3.2} is diffeomorphic to $S^3$.
Its fundamental group is trivial and its integral homology is
$(\Z,0,0,\Z)$. For negative imaginary $w$, the restricted circle action
is the standard semifree action on $S^3$ with fixed circle
$D_2^{\sigma_c}$. For positive imaginary $w$, the nonempty real form
has a free circle action isomorphic, up to circle and orientation
conventions, to the Hopf action.
\end{theorem}



\begin{proof}
\begin{enumerate}[label=\arabic*.,leftmargin=*]
\item For negative imaginary $w$ both fixed manifolds meet $D_2$ in one circle.
Their circle action is semifree, so the orbit space is a compact
orientable surface with boundary. The local quotient at a fixed circle
is the quotient of $S^1\times\C$ by scalar rotation on $\C$, hence is
a boundary collar. Over the surface interior the principal circle bundle
is trivial because the surface has boundary and $H^2=0$. Changes of
boundary trivialization extend over the attached normal discs by
rotation. Thus the three-manifold is
\[
 (B\times S^1)\ \cup_{\partial B\times S^1}\ (\partial B\times D^2).
\]
If $B$ has genus $g_B$ and $b_B$ boundary components, van Kampen or its
surface handle decomposition gives first homology
$\Z^{2g_B+b_B-1}$. Smith's theorem makes this rank zero. Since $b_B>0$,
we have $g_B=0,b_B=1$. The displayed manifold is then
$\partial(D^2\times D^2)=S^3$, with the indicated circle action. This
derives the smooth recognition from an explicit surface decomposition.

\item For positive imaginary $w$ the nonempty real manifold is free under
$\rho$. The two real tori on $t<1$ join at the single real torus over
$1$. At $0$ the odd root comparison is a torus product. At infinity
they are capped by the real neighbourhoods of $D_1$ and $D_{12}$.
These are solid tori: the primitive circle in the edge direction
collapses and the remaining real circle is their core. In the outer
real-fibre lattice $\Z\langle w,d\rangle$ the two cap meridians are
$w$ and $w-d$.

\item  Put
$B_1=A_2^{-1}A_1^{-1}A_2$ and $C=B_1^2$. The two real-arc involutions
satisfy $R_<R_m=B_1$, $R_mB_1R_m=B_1^{-1}$ and $B_1^3=I$.
Moving between the real root lifts requires
$B_1^kR_mB_1^{-k}=B_1R_m$, so $2k\equiv1\pmod3$ and $k=2$.
Thus the odd-root transport is $C$. It satisfies $R_<=CR_mC^{-1}$, and
$C(u+w)=-w+d$. The exchange across the multiplicity-four real torus
is $A_2^2$; after this odd-root transport it is
$CA_2^2C^{-1}$, whose restriction to $\langle w,d\rangle$ is
$w\mapsto-w$, $d\mapsto d$. Consequently the meridians, compared on
one boundary torus, have determinant of absolute value one. The
resulting genus-one Heegaard splitting is $S^3$: primitive meridians
intersecting once can be taken to the two coordinate circles by an
integral torus automorphism, giving the usual two solid tori of $S^3$.
The free circle $d$ then has Euler number of absolute value one, so
the displayed bundle over $S^2$ is the Hopf bundle. This proves the
fundamental-group and integral-homology assertions.
\end{enumerate}
\end{proof}
